% vim: spl=pt
\begin{exercício}{Espectro de operadores auto-adjuntos}{ex21}
    Seja \(T \in \bounded(\hilbert)\) um operador auto-adjunto. Mostre que
    \begin{enumerate}[label=(\alph*)]
        \item \(\sigma(T) \subset \mathbb{R}\) e
        \item \(\sigma_r(T) = \emptyset.\)
    \end{enumerate}
\end{exercício}
\begin{proof}[Resolução de (a)]

\end{proof}
\begin{proof}[Resolução de (b)]

\end{proof}
